% ======================================================================
% DODATAK A: OLTP POHRANJENE PROCEDURE
% ======================================================================
\chapter{OLTP pohranjene procedure}
\label{app:oltp}

Ovaj prilog sadrži implementaciju tri reprezentativne OLTP pohranjene procedure korištene u benchmarku. Odabrane su procedure koje najjasnije ilustruju arhitekturalne razlike između ispitivanih konfiguracija: najzahtjevnija transakcijska operacija (\texttt{usp\_InsertVote}), njena nativno kompajlirana varijanta za memorijsko-optimizirane sheme, te reprezentativna operacija čitanja (\texttt{usp\_GetRecentUserPosts}). Preostale tri procedure (\texttt{usp\_GetPostById}, \texttt{usp\_GetUserById}, \texttt{usp\_UpdatePostScore}) slijede identičan strukturni obrazac. Procedure za \texttt{baseline.*} i \texttt{cs.*} shemu su logički identične -- razlikuju se isključivo po prefiksu sheme.

\section{usp\_InsertVote -- disk-zasnovane konfiguracije (Baseline, CS)}

Procedura \texttt{usp\_InsertVote} predstavlja najzahtjevniju OLTP operaciju (35\% workloada) jer kombinuje dvije atomarne operacije: INSERT zapisa u tabelu \texttt{Votes} i uslovni UPDATE polja \texttt{Score} u tabeli \texttt{Posts}. Ovakav obrazac pisanja u dvije tabele unutar jednog poziva testira sposobnost storage engina da koordinira kratke, visokofrekventne transakcije. Na disk-zasnovanim konfiguracijama obe operacije zahtijevaju I/O pristup i ekskluzivno zaključavanje redova.

\begin{lstlisting}[language=SQL, caption={baseline.usp\_InsertVote -- INSERT glasa s uslovnim UPDATE score-a (35\,\% workloada)}, label=lst:a1]
CREATE OR ALTER PROCEDURE baseline.usp_InsertVote
    @Id INT, @PostId INT, @VoteTypeId TINYINT,
    @CreationDate DATE, @UserId INT
AS
BEGIN
    SET NOCOUNT ON;
    INSERT INTO baseline.Votes (Id, PostId, VoteTypeId, CreationDate, UserId)
    VALUES (@Id, @PostId, @VoteTypeId, @CreationDate, @UserId);
    -- UpVote (+1) ili DownVote (-1)
    IF @VoteTypeId = 1
        UPDATE baseline.Posts SET Score = Score + 1 WHERE Id = @PostId;
    ELSE IF @VoteTypeId = 2
        UPDATE baseline.Posts SET Score = Score - 1 WHERE Id = @PostId;
END;
\end{lstlisting}

\section{im.usp\_InsertVote -- memorijsko-optimizirane konfiguracije (IM, HTAP)}

Za memorijsko-optimizirane sheme (\texttt{im.*}, \texttt{htap.*}) ista procedura deklarira se s klauzulama \texttt{WITH NATIVE\_COMPILATION, SCHEMABINDING}. Nativna kompilacija prevodi T-SQL u strojni kod u trenutku kreiranja procedure, eliminirajući T-SQL interpretacijski overhead pri svakom pozivu. \texttt{BEGIN ATOMIC} blok garantira atomarnost bez potrebe za eksplicitnim \texttt{BEGIN TRANSACTION}/\texttt{COMMIT}, a \texttt{TRANSACTION ISOLATION LEVEL = SNAPSHOT} osigurava optimistički MVCC bez zaključavanja čitatelja. Logika procedure identična je disk-zasnovanoj varijanti -- razlika je isključivo u načinu izvršavanja.

\begin{lstlisting}[language=SQL, caption={im.usp\_InsertVote -- ista logika, nativno kompajlirana s ATOMIC blokom}, label=lst:a2]
IF OBJECT_ID('im.usp_InsertVote', 'P') IS NOT NULL
    DROP PROCEDURE im.usp_InsertVote;
GO
CREATE PROCEDURE im.usp_InsertVote
    @Id INT, @PostId INT, @VoteTypeId TINYINT,
    @CreationDate DATE, @UserId INT
WITH NATIVE_COMPILATION, SCHEMABINDING
AS BEGIN ATOMIC WITH
    (TRANSACTION ISOLATION LEVEL = SNAPSHOT, LANGUAGE = N'English')
    INSERT INTO im.Votes (Id, PostId, VoteTypeId, CreationDate, UserId)
    VALUES (@Id, @PostId, @VoteTypeId, @CreationDate, @UserId);
    IF @VoteTypeId = 1
        UPDATE im.Posts SET Score = Score + 1 WHERE Id = @PostId;
    ELSE IF @VoteTypeId = 2
        UPDATE im.Posts SET Score = Score - 1 WHERE Id = @PostId;
END;
\end{lstlisting}

\section{usp\_GetRecentUserPosts -- range scan s TOP/ORDER BY}

Procedura \texttt{usp\_GetRecentUserPosts} (10\% workloada) predstavlja tipičnu read operaciju na web platformi: dohvaćanje zadnjih 10 postova određenog korisnika sortirano po datumu. Za razliku od point lookup-a, ova operacija izvodi \textit{range scan} -- filtrira po \texttt{OwnerUserId} (ne-PK stupac) i sortira rezultat, što zahtijeva nonclustered indeks na disk-zasnovanim konfiguracijama, a na memorijsko-optimiziranim tabelama podrazumijeva sekvencijalni prolaz kroz hash indeks.

\begin{lstlisting}[language=SQL, caption={baseline.usp\_GetRecentUserPosts -- range scan s TOP 10 i ORDER BY (10\,\% workloada)}, label=lst:a3]
CREATE OR ALTER PROCEDURE baseline.usp_GetRecentUserPosts
    @UserId INT
AS
BEGIN
    SET NOCOUNT ON;
    SELECT TOP 10 Id, PostTypeId, CreationDate, Score,
                  Title, AnswerCount, CommentCount
    FROM baseline.Posts
    WHERE OwnerUserId = @UserId
    ORDER BY CreationDate DESC;
END;
\end{lstlisting}
